\section{Results}
\label{sec:results}

We investigate the dependence of outer rotation curve slopes ($\mathrm{slope\_tail}$) on
environmental density using a sample of $N=79$ galaxies from the SPARC catalogue
\citep{Lelli2016}.

% -------------------------------------------------------------------
\subsection{Global correlation}
\label{sec:results:global}

We find a statistically significant anti-correlation between $\mathrm{slope\_tail}$
and the environmental proxy over the full sample:

\begin{equation}
  \rho = -0.365,\quad p = 9.3 \times 10^{-4}.
\end{equation}

Galaxies residing in denser environments tend to exhibit steeper (more negative)
outer rotation-curve slopes, suggesting that the large-scale tidal field modulates
the outer kinematics.

% -------------------------------------------------------------------
\subsection{Mass-dependent regime}
\label{sec:results:mass}

Splitting the sample by baryonic mass reveals a clear transition in the
environmental signal (Figure~\ref{fig:main}, left panel):

\begin{itemize}
  \item \emph{Low-mass} galaxies ($\log M_{\rm bar} < 10$): no measurable
        correlation with environment ($\rho \approx 0$, $p \gg 0.05$).
  \item \emph{High-mass} galaxies ($\log M_{\rm bar} > 10.5$): strong
        anti-correlation.
\end{itemize}

For the high-mass subsample ($N = 47$):

\begin{equation}
  \rho = -0.42,\quad p = 0.003.
\end{equation}

The transition is not gradual; the signal emerges abruptly above a critical
baryonic mass scale $\log M_{\rm bar} \simeq 10.5$.

% -------------------------------------------------------------------
\subsection{Regression analysis}
\label{sec:results:regression}

We perform an OLS regression with heteroscedasticity-robust standard errors
(HC3 estimator) restricted to the high-mass subsample
(Figure~\ref{fig:main}, right panel):

\begin{equation}
  \mathrm{slope\_tail} = \beta_0 + \beta_{\rm env}\,\mathrm{env\_proxy}.
\end{equation}

The best-fit environmental coefficient is

\begin{equation}
  \beta_{\rm env} = -0.14 \pm 0.05,\quad p = 0.006,
\end{equation}

confirming that the environmental term is statistically required and cannot
be attributed to a chance fluctuation.

% -------------------------------------------------------------------
\subsection{Failure of continuous interaction models}
\label{sec:results:interaction}

We additionally test a model that includes a continuous $\mathrm{env} \times M$
interaction term across the full mass range.  This model fails to produce a
significant interaction coefficient, indicating that the coupling between
environment and outer slope is \emph{not} a smoothly varying function of mass,
but rather switches on above a threshold.

% -------------------------------------------------------------------
\subsection{Physical interpretation}
\label{sec:results:interpretation}

The results are consistent with a two-regime dynamical picture:

\begin{itemize}
  \item \emph{Internal-dominated regime} (low mass): rotation curves are
        governed by the internal baryonic mass distribution; environmental
        effects are negligible.
  \item \emph{Environment-coupled regime} (high mass): the outer rotation
        curve is measurably modulated by the large-scale tidal environment.
\end{itemize}

The abrupt onset of the signal at a critical mass scale suggests a structural
transition in the way angular momentum and tidal torques couple, rather than a
continuous environmental feedback.

% -------------------------------------------------------------------
\subsection{Summary}
\label{sec:results:summary}

\begin{table}
  \centering
  \caption{Summary of statistical results.}
  \label{tab:results}
  \begin{tabular}{lcccc}
    \hline
    Subsample & $N$ & $\rho$ & $p$-value & $\beta_{\rm env}$ \\
    \hline
    Full sample   & 79 & $-0.365$ & $9.3\times10^{-4}$ & --- \\
    Low mass      & 32 & $\approx 0$ & $> 0.05$ & --- \\
    High mass     & 47 & $-0.42$  & $0.003$            & $-0.14\pm0.05$ \\
    \hline
  \end{tabular}
\end{table}

The SCM framework reveals a previously hidden dynamical structure in galaxy
rotation curves: environmental modulation is real and statistically robust, but
it is regime-dependent rather than universal.
